\section{Introduction}

Despite the rapid evolution of audio-visual large language models (AV-LLMs), most existing systems continue to rely on RGB video and monaural audio \cite{xu2025qwen2, Qwen3-Omni, tang2025video, cheng2024videollama}, a design choice that fundamentally constrains their ability to perceive and reason about the three-dimensional (3D) physical world. Although 3D grounding has recently attracted significant attention, current research remains fragmented, with most approaches addressing spatial cues from vision or audio in isolation.
In the visual domain, recent advances extend a 2D visual LLM by incorporating RGB-D inputs and chain-of-thought reasoning, leading to improved 3D grounding and spatial relationship understanding \cite{wang2025n3d}. In parallel, auditory approaches leverage binaural spatial encoders \citep{zheng2024bat} or intensity-vector representations \citep{tang2024can} to enable relative sound source localization. However, a unified paradigm for comprehensive 3D audio-visual scene understanding remains largely unexplored.
Early multimodal efforts, such as Hear You Are \citep{ryu2025hear}, 
pair spatial audio with panoramic RGB image and assume a single active source per scene, which precludes evaluating overlap-robust localization and depth-aware 3D grounding. 
Similarly, while SAVVY \citep{chen2025savvy} integrates RGB-D perception with multi-channel audio, it relies on a cascaded pipeline that depends on traditional signal processing for sound source localization. 
This modular design hinders end-to-end learning and prevents the AV-LLM from performing fully integrated spatial reasoning, underscoring the need for a more unified, learned approach.

In this paper, we introduce JAEGER, a \textbf{j}oint 3D \textbf{a}udio-visual r\textbf{e}asoning and \textbf{g}rounding method in simulated physical \textbf{e}nvi\textbf{r}onments, an end-to-end framework that extends 2D audio-visual large language models (AV-LLMs) to 3D settings by explicitly modeling visual depth and multi-channel spatial audio. Initialized from Qwen2.5-Omni \citep{xu2025qwen2} and efficiently adapted via low-rank adaptation (LoRA) \cite{hu2022lora}, JAEGER integrates an RGB-D visual stream augmented with depth-projected 3D positional encodings together with a dual-path audio stream, in which semantic content extracted from the omnidirectional first-order ambisonics (FOA) channel is disentangled from spatial directional cues.
To further enhance azimuth perception in reverberant environments and under overlapping sound sources, we propose a novel neural intensity vector (Neural IV) approach, which is a learnable FOA-based representation that encodes robust directional information for improved sound localization and cross-modal alignment.

To support instruction tuning and systematic evaluation, we construct SpatialSceneQA, a benchmark comprising 61k high-fidelity RGB-D scenes paired with spatial audio and fine-grained 3D annotations. The dataset covers a diverse set of tasks, including single- and overlapping-source direction-of-arrival (DoA) estimation, 3D visual grounding of sound-emitting objects, and multi-speaker audio-visual matching. Experimental results show that JAEGER achieves a median angular error (MAE) of 2.21° for single sources and 13.13° under overlapping conditions, while showing strong generalization across varied source configurations. Leveraging explicit depth cues, the model attains a 3D intersection over union (IoU) of 0.32 with a median localization error of 0.16 meter (m). When FOA-based spatial cues are jointly modeled with RGB-D perception, JAEGER further achieves 99.2\% accuracy on joint audio-visual reasoning in simulated multi-speaker physical environments.

In summary, our main contributions are as follows:
\vspace{-0.2cm}
\begin{itemize}[leftmargin=*]
    \setlength\itemsep{-0.3em}
    
    \item We introduce joint grounding and reasoning with spatial audio and RGB-D geometry, which exposes a fundamental modality gap for current AV-LLMs even after fine-tuning. We adapt an AV-LLM with depth-aware visual encoding and FOA spatial cues for end-to-end DoA estimation, 3D box grounding, and multi-speaker matching.
    \item We present \textsc{SpatialSceneQA}, a 61k-sample dataset that pairs 3D RGB-D images with 3D 4-channel FOA audio and dense object-level spatial annotations. To the best of our knowledge, it is the first spatial audio-visual benchmark with degree-level azimuth and elevation supervision, enabling precise localization, 3D grounding, and multi-speaker matching under source overlap.
    \item We propose the \textit{Neural Intensity Vector}, a learnable FOA spatial representation that replaces STFT-based intensity features with a neural encoder, yielding more stable azimuth cues under reverberation and overlapping sources and improving cross-scenario generalization.
\end{itemize}
